\subsection*{Real symmetric spectra}

\begin{problem}[Symmetric orthogonal matrices]\label{ex:09:symmetric-orthogonal}
Prove that a real symmetric orthogonal matrix has eigenvalues only
\(1\) and \(-1\). If it is also positive definite, prove it is \(I\).
\end{problem}
\begin{solution}
The real spectral theorem gives an eigenbasis. Since \(A^\top A=I\)
and \(A^\top=A\), we have \(A^2=I\). For an eigenvector \(v\ne0\),
\(v=A^2v=\lambda^2v\), so \(\lambda^2=1\).
If \(A\) is PD, every eigenvalue is positive, hence \(1\).
The real spectral theorem then gives \(A=QI Q^\top=I\).
\end{solution}

\begin{problem}[Independence and orthogonality with parameters]\label{ex:09:parameters}
Let \(A=\begin{pmatrix}a&1\\ab&b\end{pmatrix}\), where real \(a,b\ne0\).
\begin{enumerate}
\item Find the eigenvalues and normalized eigenvectors.
\item Show that eigenvectors for the two roots are independent unless
\(a+b=0\).
\item Show that these directions are orthogonal exactly when \(ab=1\).
\end{enumerate}
\end{problem}
\begin{solution}
The polynomial is \(t^2-(a+b)t=t(t-a-b)\).
Vectors \(v_0=(1,-a)^\top\) and \(v_{a+b}=(1,b)^\top\) satisfy
the eigenvector equations. Their unit versions are
\(v_0/\sqrt{1+a^2}\) and \(v_{a+b}/\sqrt{1+b^2}\), with either sign.
Their column determinant is \(a+b\); if \(a+b\ne0\), they are independent.
If \(a+b=0\), they coincide and the single eigenspace for zero is
\(\operatorname{span}(1,-a)^\top\): \(A\ne0\) has rank one.
Their dot product before normalization is \(1-ab\), proving (c).
The condition \(ab=1\) is exactly the symmetry condition for \(A\),
consistent with the symmetric eigenvector theorem.
\end{solution}

\begin{problem}[A repeated eigenvalue]\label{ex:09:multiplicity}
Let \(\lambda\) have algebraic multiplicity \(k\geqslant2\) in a real
symmetric matrix.
\begin{enumerate}
\item Prove its eigenspace has an orthonormal basis of \(k\) vectors.
\item Prove there are infinitely many such orthonormal bases.
\item Prove there cannot be more than \(k\) independent eigenvectors
for \(\lambda\), and conclude the eigenspace has dimension \(k\).
\end{enumerate}
\end{problem}
\begin{solution}
The spectral theorem diagonalizes the matrix. The \(k\) columns belonging
to \(\lambda\) form an orthonormal basis of that eigenspace.
If the first two are \(q_1,q_2\), replace them by
\(q_1\cos t+q_2\sin t\) and \(-q_1\sin t+q_2\cos t\).
They remain orthonormal and in the same eigenspace, giving infinitely
many bases as \(t\) varies.
The earlier geometric-multiplicity bound gives dimension at most \(k\),
while the \(k\) constructed vectors give dimension at least \(k\).
\end{solution}

\begin{problem}[An orthonormal matrix spanning a kernel]\label{ex:09:null-basis}
Let \(A\in\R^{m\times n}\) have rank \(r\).
\begin{enumerate}
\item Use the eigenspaces of \(A^\top A\) to construct
\(S\in\R^{n\times(n-r)}\) with \(AS=0\) and \(S^\top S=I_{n-r}\).
\item If \(AT=0\), prove \(\operatorname{rank}T\leqslant n-r\).
\end{enumerate}
\end{problem}
\begin{solution}
The Gram kernel identity gives
\(\ker(A^\top A)=\ker A\), of dimension \(n-r\).
Choose orthonormal eigenvectors for its zero eigenvalue and place them
in \(S\). They give both stated identities.
Every column of \(T\) lies in \(\ker A\), so its column-space dimension
is at most \(n-r\).
When \(r=n\), \(S\) has no columns and the identities use the empty
zero-dimensional identity.
\end{solution}

\begin{problem}[*Simultaneous orthogonal diagonalization]\label{ex:09:simultaneous}
For real symmetric \(A,B\) of the same order, prove that one orthogonal
matrix diagonalizes both if and only if \(AB=BA\).
\end{problem}
\begin{solution}
If \(A=Q\Lambda Q^\top\), \(B=QM Q^\top\) with diagonal \(\Lambda,M\),
their products commute because \(\Lambda M=M\Lambda\).
Conversely, commutation implies \(B(E_\lambda(A))\subseteq E_\lambda(A)\).
On each such eigenspace, the restriction of \(B\) is self-adjoint.
The spectral theorem gives an orthonormal basis there of \(B\)-eigenvectors,
which are also \(A\)-eigenvectors. Different \(A\)-eigenspaces are
orthogonal, so the union of these bases is an orthonormal basis of the
whole space. Its column matrix diagonalizes both.
\end{solution}

\begin{problem}[*Craig--Sakamoto identity]\label{ex:09:craig}
For real symmetric \(A,B\) of the same order, prove
\[
 AB=0\quad\Longleftrightarrow\quad
 \det(I-sA)\det(I-tB)=\det(I-sA-tB)
 \quad\hbox{for every real }s,t.
\]
\end{problem}
\begin{solution}
If \(AB=0\), then \((I-sA)(I-tB)=I-sA-tB\); take determinants.

For the converse, work for all \(s\) in a sufficiently small interval
about zero, where \(C(s)=(I-sA)^{-1}\) exists.
The assumed identity, divided by \(\det(I-sA)\), gives
\(\det(I-tB)=\det(I-tC(s)B)\) for every \(t\).
Consequently \(B\) and \(C(s)B\) have the same characteristic polynomial,
and
\[
 h(s):=\tr\{C(s)BC(s)B\}=\tr B^2 .
\]
The trace-of-square conclusion follows directly from the coefficients
of the determinant, without any hypothesis on its roots:
the first coefficient gives \(\tr M\), and the second gives
\(\{(\tr M)^2-\tr(M^2)\}/2\), by summing the two-by-two principal minors.
For order one, the first coefficient already gives \(M\) itself.

The entries of \(C(s)\) are differentiable rational functions near zero.
Differentiating \((I-sA)C(s)=I\) gives
\(C'(0)=A\), \(C''(0)=2A^2\).
Twice applying the product rule to \(h\), and cyclically moving factors
inside the trace, yields
\[
 0=h''(0)=4\tr(A^2B^2)+2\tr(ABAB).
\]
Put \(M=AB\). Symmetry gives \(\tr(A^2B^2)=\tr(MM^\top)=\norm M_F^2\),
where \(\norm M_F^2=\sum_{i,j}m_{ij}^2\).
Also
\[
 \tr(M^2)+\norm M_F^2=\tfrac12\norm{M+M^\top}_F^2.
\]
Therefore
\[
 0=\tr(M^2)+2\norm M_F^2
   =\tfrac12\norm{M+M^\top}_F^2+\norm M_F^2.
\]
Both terms are nonnegative, so \(M=0\).
The proof does not assume \(\tr(M^2)\geqslant0\) for a nonsymmetric
matrix; that assertion would be false.
\end{solution}

\begin{problem}[The norm of a matrix acting on a vector]\label{ex:09:norm-bound}
For real \(A\in\R^{m\times n}\), let \(\mu\) be the largest eigenvalue
of \(A^\top A\). Prove \(\norm{Ax}\leqslant\sqrt\mu\,\norm x\) for
every \(x\).
\end{problem}
\begin{solution}
The Gram matrix is PSD, so \(\mu\geqslant0\).
For \(x\ne0\), the Rayleigh bound gives
\(\norm{Ax}^2=x^\top A^\top Ax\leqslant\mu x^\top x\).
Taking nonnegative square roots proves the result.
For \(x=0\), both sides vanish. Equality for \(x\ne0\) occurs precisely
in the top eigenspace of \(A^\top A\).
\end{solution}

\subsection*{Positive forms: basic tests and constructions}

\begin{problem}[Quadratic-form uniqueness]\label{ex:09:quadratic-uniqueness}
\begin{enumerate}
\item For real square \(A\), show
\(x^\top Ax=x^\top(A+A^\top)x/2\).
\item Prove \(x^\top Ax=0\) for all \(x\) exactly when \(A\) is skew-symmetric.
\item Deduce that two symmetric matrices with the same quadratic form
are equal.
\end{enumerate}
\end{problem}
\begin{solution}
Transposing a scalar gives \(x^\top Ax=x^\top A^\top x\), proving (a).
If \(A^\top=-A\), that scalar equals its negative and is zero.
Conversely, setting \(x=e_i\) gives \(a_{ii}=0\);
setting \(x=e_i+e_j\) gives \(a_{ij}+a_{ji}=0\).
Thus \(A^\top=-A\).
If \(A,B\) are symmetric and have the same form, \(A-B\) is both
symmetric and skew-symmetric, hence zero.
\end{solution}

\begin{problem}[Represent and classify three forms]\label{ex:09:forms}
Write the following as symmetric matrix forms and classify their signs:
\[
 x^2+y^2-xy,\qquad
 4x^2+5y^2+z^2+2xy+2yz,\qquad
 -x^2+2y^2+xy.
\]
\end{problem}
\begin{solution}
Their matrices are
\[
 A=\begin{pmatrix}1&-1/2\\-1/2&1\end{pmatrix},\quad
 B=\begin{pmatrix}4&1&0\\1&5&1\\0&1&1\end{pmatrix},\quad
 C=\begin{pmatrix}-1&1/2\\1/2&2\end{pmatrix}.
\]
The first form is \((x-y/2)^2+3y^2/4\), positive for nonzero \((x,y)\).
The second is
\[
 4(x+y/4)^2+(z+y)^2+\frac{15}{4}y^2,
\]
also positive for every nonzero vector.
The third takes values \(-1\) at \((1,0)\) and \(2\) at \((0,1)\),
so it is indefinite.
\end{solution}

\begin{problem}[Length preservation]\label{ex:09:isometry}
For a fixed real square matrix \(Q\), prove that the following are equivalent:
\(Q^\top Q=I\); \(\norm{Qx}=\norm x\) for every \(x\);
\((Qx)^\top(Qy)=x^\top y\) for every \(x,y\).
Does equality of lengths for just one vector imply orthogonality?
\end{problem}
\begin{solution}
\(Q^\top Q=I\) gives the inner-product identity, which gives norm
preservation on taking \(y=x\).
Conversely, norm preservation makes
\(x^\top(Q^\top Q-I)x=0\) for all \(x\).
The matrix is symmetric, so the preceding quadratic-form uniqueness
result gives \(Q^\top Q-I=0\).
A single vector is insufficient: \(Q=\operatorname{diag}(1,2)\)
preserves the length of \(e_1\) but is not orthogonal.
The matrix must be fixed and the condition must hold for all vectors.
\end{solution}

\begin{problem}[Signs and diagonal matrices]\label{ex:09:signs}
\begin{enumerate}
\item Prove \(A\succ0\) iff \(-A\prec0\), and \(A\succeq0\) iff
\(-A\preceq0\).
\item Can a matrix be both positive and negative definite? Both positive
and negative semidefinite?
\item For diagonal \(D=\operatorname{diag}(d_i)\), characterize PD and PSD.
\end{enumerate}
\end{problem}
\begin{solution}
Multiplication by \(-1\) reverses the sign of every quadratic value,
proving (a).
A positive-dimensional matrix cannot have both strictly positive and
strictly negative quadratic values on a fixed nonzero vector.
Both weak signs hold exactly when its quadratic form vanishes, which
for symmetric \(A\) means \(A=0\).
For a diagonal matrix \(x^\top Dx=\sum_i d_ix_i^2\).
If all \(d_i>0\) it is PD; if all \(d_i\geqslant0\) it is PSD.
Testing \(x=e_i\) proves the necessity of those respective conditions.
\end{solution}

\begin{problem}[*Entrywise bounds]\label{ex:09:entries}
Let \(A=(a_{ij})\succ0\).
\begin{enumerate}
\item Prove \(a_{ii}>0\).
\item For \(i\ne j\), prove \(a_{ij}^2<a_{ii}a_{jj}\).
\item Prove a largest absolute entry of \(A\) is a diagonal entry.
\item State and prove the corresponding PSD conclusions.
\end{enumerate}
\end{problem}
\begin{solution}
Testing \(e_i\) gives (a).
For \(x=te_i+e_j\), positivity gives
\(a_{ii}t^2+2a_{ij}t+a_{jj}>0\) for every real \(t\).
Its minimum at \(t=-a_{ij}/a_{ii}\) is
\(a_{jj}-a_{ij}^2/a_{ii}>0\), proving (b).
Therefore
\(|a_{ij}|<\sqrt{a_{ii}a_{jj}}\leqslant\max_k a_{kk}\).
For PSD matrices replace the strict signs by weak signs.
If \(a_{ii}=0\), the expression \(2a_{ij}t+a_{jj}\) must be nonnegative
for every \(t\), which forces \(a_{ij}=0\). If \(a_{ii}>0\), minimize
as before to obtain \(a_{ij}^2\leqslant a_{ii}a_{jj}\).
These prove the diagonal maximum statement also in the PSD case.
\end{solution}

\begin{problem}[Sums, trace, and determinant]\label{ex:09:psd-sums}
Let \(A\succeq0\).
\begin{enumerate}
\item Prove \(\tr A\geqslant0\), with equality exactly when \(A=0\).
\item Prove \(\det A\geqslant0\), with equality exactly when \(A\) is singular.
\item If \(B\succeq0\), prove \(A+B\succeq0\); if \(B\succ0\),
prove \(A+B\succ0\).
\end{enumerate}
\end{problem}
\begin{solution}
The eigenvalues \(\lambda_i\) are nonnegative. Their sum is nonnegative
and vanishes exactly when every one is zero; spectral decomposition
then gives \(A=0\).
Their product is nonnegative and vanishes exactly when an eigenvalue is
zero, equivalently when \(A\) is singular.
Finally \(x^\top(A+B)x=x^\top Ax+x^\top Bx\).
The terms are nonnegative, and the second is strictly positive for
\(x\ne0\) if \(B\succ0\), proving (c).
\end{solution}

\begin{problem}[When a Gram matrix is zero]\label{ex:09:zero-gram}
For real \(A\in\R^{m\times n}\), prove \(A^\top A,AA^\top\succeq0\) and
\[
 A=0\quad\Longleftrightarrow\quad A^\top A=0
 \quad\Longleftrightarrow\quad\tr(A^\top A)=0.
\]
\end{problem}
\begin{solution}
Their quadratic forms are \(\norm{Ax}^2\) and \(\norm{A^\top y}^2\).
Also
\(\tr(A^\top A)=\sum_{i,j}a_{ij}^2\).
This sum vanishes precisely when all entries vanish.
The forward implications to a zero Gram matrix and zero trace are
immediate, and the sum-of-squares formula proves the reverse implication.
\end{solution}

\begin{problem}[Positive determinant and trace are insufficient]\label{ex:09:det-trace}
Prove that a PD matrix has positive determinant and positive trace.
Show that these two scalar conditions together do not imply positive
definiteness, even for a real symmetric matrix.
\end{problem}
\begin{solution}
All eigenvalues of a PD matrix are positive, so their product and sum
are positive. Conversely
\(A=\operatorname{diag}(-1,-1,3)\) has determinant \(3>0\) and
trace \(1>0\), but \(e_1^\top Ae_1=-1<0\); it is indefinite.
\end{solution}

\begin{problem}[Inverses and powers preserve positivity]\label{ex:09:inverse-powers}
\begin{enumerate}
\item Prove a PD matrix is invertible, and a PSD matrix is PD exactly
when it is invertible.
\item For symmetric invertible \(A\), prove \(A\succ0\) iff \(A^{-1}\succ0\).
\item If \(A\) is PD or PSD, prove the same respective property for
\(A^k\), \(k=1,2,\ldots\).
\end{enumerate}
\end{problem}
\begin{solution}
In the spectral representation, invertibility means that every eigenvalue
is nonzero. A nonnegative eigenvalue list is strictly positive exactly
when none is zero, proving (a).
The eigenvalues of \(A^{-1}\) are \(1/\lambda_i\), which are positive
exactly when all \(\lambda_i\) are positive; the inverse is symmetric.
The powers have the same orthonormal eigenvectors and eigenvalues
\(\lambda_i^k\), which retain positivity or nonnegativity.
\end{solution}

\begin{problem}[Complete a matrix square]\label{ex:09:complete-square}
Let \(A,B\succ0\), \(C=(A^{-1}+B^{-1})^{-1}\), and \(D=A+B\).
\begin{enumerate}
\item Prove
\[
 D^{-1}=A^{-1}-A^{-1}CA^{-1}
 =B^{-1}-B^{-1}CB^{-1}
 =A^{-1}CB^{-1}=B^{-1}CA^{-1}.
\]
\item Find \(\mu\) and an \(x\)-independent quantity \(Q(a,b)\) such that
\[
 (x-a)^\top A^{-1}(x-a)+(x-b)^\top B^{-1}(x-b)
 =(x-\mu)^\top C^{-1}(x-\mu)+Q(a,b).
\]
\item Prove \(Q(a,b)=(a-b)^\top D^{-1}(a-b)\).
\end{enumerate}
\end{problem}
\begin{solution}
All inverses exist: sums of PD matrices are PD, as are their inverses.
Since \(A^{-1}DB^{-1}=B^{-1}+A^{-1}=C^{-1}\),
inverting gives \(C=BD^{-1}A\).
Interchanging \(A,B\) also gives \(C=AD^{-1}B\).
Multiplying these identities gives
\(A^{-1}CB^{-1}=B^{-1}CA^{-1}=D^{-1}\).
Furthermore \(C(A^{-1}+B^{-1})=I\), so
\[
 A^{-1}-A^{-1}CA^{-1}
 =A^{-1}(I-CA^{-1})=A^{-1}CB^{-1}=D^{-1};
\]
the other formula follows by interchange.

Set \(h=A^{-1}a+B^{-1}b\) and \(\mu=Ch\).
Expansion of both sides of (b) gives the same quadratic term
\(x^\top C^{-1}x\) and linear term \(-2x^\top h\), provided
\[
 Q(a,b)=a^\top A^{-1}a+b^\top B^{-1}b-h^\top Ch.
\]
Using (a), the coefficients of the pure \(a\) and pure \(b\) terms
are \(D^{-1}\), while the two cross terms are
\(-a^\top D^{-1}b-b^\top D^{-1}a\).
Thus \(Q(a,b)=(a-b)^\top D^{-1}(a-b)\), as claimed.
\end{solution}

\begin{problem}[*A shift and a gap in the spectrum]\label{ex:09:kato}
\begin{enumerate}
\item For symmetric \(A\), show that some real \(k\) makes \(kI+A\) PD.
\item If symmetric \(A\) has no eigenvalue in \([a,b]\), where \(a\leqslant b\),
prove \((A-aI)(A-bI)\succ0\).
\end{enumerate}
\end{problem}
\begin{solution}
Choose \(k>-\lambda_{\min}(A)\); the shifted eigenvalues \(k+\lambda_i\)
are all positive.
For (b), the two factors commute and are simultaneously diagonal in an
eigenbasis of \(A\). Their product has eigenvalues
\((\lambda_i-a)(\lambda_i-b)\).
Since \(\lambda_i<a\) or \(\lambda_i>b\), both factors have the same
strict sign. Every product is positive, proving the claim, also when
\(a=b\).
\end{solution}

\begin{problem}[A sum of two rank-one matrices]\label{ex:09:rank-two-plus}
Let \(a,b\in\R^n\) be independent and \(A=aa^\top+bb^\top\).
Prove \(A\succeq0\) with rank two. Find its two positive eigenvalues
through a two-by-two equation, and express the corresponding eigenvectors
as combinations of \(a,b\).
\end{problem}
\begin{solution}
Let \(X=(a\ b)\). Then \(A=XX^\top\succeq0\), and the Gram rank identity
gives \(\operatorname{rank}A=\operatorname{rank}X=2\).
Its image is \(\operatorname{span}(a,b)\). A nonzero eigenvalue has
eigenvector \(x=\lambda^{-1}Ax\) in that image, so write
\(x=\theta_1a+\theta_2b\).
Independence makes \(Ax=\lambda x\) equivalent to
\[
 \begin{pmatrix}\lambda-a^\top a&-a^\top b\\
                 -a^\top b&\lambda-b^\top b\end{pmatrix}
 \begin{pmatrix}\theta_1\\\theta_2\end{pmatrix}=0.
\]
The Gram matrix \(X^\top X\) is PD, so its two eigenvalues are positive;
they are the roots of
\[
 \lambda^2-(\norm a^2+\norm b^2)\lambda
 +\norm a^2\norm b^2-(a^\top b)^2=0.
\]
Every nonzero solution \(\theta\) produces a nonzero eigenvector \(X\theta\).
All remaining \(n-2\) eigenvalues are zero.
\end{solution}

\begin{problem}[A difference of two rank-one matrices]\label{ex:09:rank-two-minus}
Let \(a,b\in\R^n\) be independent and \(A=aa^\top-bb^\top\).
Prove its rank is two and its two nonzero eigenvalues have opposite signs.
Derive the equation
\[
 \begin{pmatrix}\lambda-a^\top a&a^\top b\\
                 -a^\top b&\lambda+b^\top b\end{pmatrix}
 \begin{pmatrix}\theta_1\\\theta_2\end{pmatrix}=0
\]
and show the corresponding eigenvectors are \(x=\theta_1a-\theta_2b\).
\end{problem}
\begin{solution}
The image lies in \(\operatorname{span}(a,b)\).
In the basis \(a,-b\) of this plane, the restriction has matrix
\[
 K=\begin{pmatrix}\norm a^2&-a^\top b\\
                  a^\top b&-\norm b^2\end{pmatrix}.
\]
Its determinant is
\((a^\top b)^2-\norm a^2\norm b^2<0\), by strict Schwarz for independent
vectors. Thus its restriction is invertible, so the rank is two.
Its characteristic polynomial has real coefficients and negative
constant term; its discriminant is the squared trace minus four times
that negative determinant, hence positive. The two real roots have
negative product and therefore opposite signs.
The displayed equation is \((\lambda I-K)\theta=0\), and
\(x=\theta_1a-\theta_2b\) translates it back to the original space.
Thus the form is indefinite.
\end{solution}

\begin{problem}[Factors of a PSD matrix]\label{ex:09:psd-factor}
Let \(A\succeq0\) have order \(n\) and rank \(r\).
\begin{enumerate}
\item Construct \(T\in\R^{n\times r}\) and positive diagonal
\(\Lambda\in\R^{r\times r}\) with \(A=TT^\top\), \(T^\top T=\Lambda\).
\item For each \(r\leqslant m\leqslant n\), construct an \(n\times m\)
matrix \(B\) with \(A=BB^\top\).
\end{enumerate}
\end{problem}
\begin{solution}
Take the orthonormal eigenvectors \(q_1,\ldots,q_r\) for the positive
eigenvalues and put
\(T=(\sqrt{\lambda_1}q_1\ \cdots\ \sqrt{\lambda_r}q_r)\).
Then \(TT^\top=\sum_{j=1}^r\lambda_jq_jq_j^\top=A\), and
\(T^\top T=\operatorname{diag}(\lambda_1,\ldots,\lambda_r)\).
If \(A\) is PD, \(r=n\) and \(T\) is square and invertible.
For (b), append \(m-r\) zero columns to \(T\).
For \(A=0\), use the empty matrix when \(r=m=0\), or a zero matrix
of the requested positive width.
\end{solution}

\begin{problem}[Cholesky factorization in the definite case]\label{ex:09:cholesky}
Prove that every real \(A\succ0\) can be written \(A=LL^\top\), where
\(L\) is lower triangular with positive diagonal entries.
\end{problem}
\begin{solution}
Induct on the order. For \(n=1\), take \(L=[\sqrt{a_{11}}]\).
Partition
\(A=\begin{pmatrix}a&b^\top\\b&D\end{pmatrix}\), where \(a>0\).
For \(y\ne0\), test the nonzero vector
\((-b^\top y/a,y^\top)^\top\); positivity gives
\(y^\top(D-bb^\top/a)y>0\).
Thus \(S=D-bb^\top/a\) is PD.
By induction \(S=L_2L_2^\top\), with \(L_2\) lower triangular and
positive diagonal. Set
\[
 L=\begin{pmatrix}\sqrt a&0\\b/\sqrt a&L_2\end{pmatrix}.
\]
Multiplication gives \(LL^\top=A\), and every diagonal entry is positive.
This proves precisely the positive-definite case.
\end{solution}

\begin{problem}[*The unique PSD square root]\label{ex:09:square-root}
\begin{enumerate}
\item Construct a PSD \(B\) with \(B^2=A\) when \(A\succeq0\).
\item Prove this PSD \(B\) is unique, although unrestricted square roots
need not be unique.
\item If \(A\succ0\), prove
\((A^{-1})^{1/2}=(A^{1/2})^{-1}\).
\end{enumerate}
\end{problem}
\begin{solution}
An orthonormal eigenbasis gives
\(B=Q\operatorname{diag}(\sqrt{\lambda_i})Q^\top\), proving existence
and providing a calculation method.
If another PSD \(C\) satisfies \(C^2=A\), then \(CA=AC\).
Therefore each eigenspace \(E_\lambda(A)\) is invariant under \(C\).
On this space \(C\) is self-adjoint and PSD, so it has an orthonormal
eigenbasis with eigenvalues \(\mu\geqslant0\).
The identity \(C^2=A\) gives \(\mu^2=\lambda\), hence
\(\mu=\sqrt\lambda\).
Thus \(C\) acts as \(\sqrt\lambda I\) on every \(E_\lambda(A)\), exactly
as \(B\) does; the eigenspaces span the whole space, so \(C=B\).
For example \(I\) has the unrestricted square roots \(I\) and \(-I\),
showing the PSD restriction matters.
When \(A\succ0\), both matrices in (c) have eigenvectors \(q_i\)
and eigenvalues \(1/\sqrt{\lambda_i}\), proving their equality.
\end{solution}

\begin{problem}[Matrix quadratic forms]\label{ex:09:matrix-forms}
Let \(A\succeq0\) be \(n\times n\) and \(B\in\R^{n\times m}\).
\begin{enumerate}
\item Prove \(B^\top AB\succeq0\) and
\(\operatorname{rank}(B^\top AB)=\operatorname{rank}(AB)\).
\item Prove \(B^\top AB=0\) iff \(AB=0\).
\item If \(A\succ0\), prove
\(\operatorname{rank}(B^\top AB)=\operatorname{rank}B\);
deduce the conditions for \(B^\top AB\succ0\) and for \(B^\top AB=0\).
\item Even if \(B\) has full column rank \(m<n\), does
\(B^\top AB\succ0\) imply \(A\succ0\)?
\end{enumerate}
\end{problem}
\begin{solution}
Set \(C=A^{1/2}B\). Then \(B^\top AB=C^\top C\succeq0\), with
kernel \(\ker C\).
Spectral decomposition gives \(\ker A^{1/2}=\ker A\).
Therefore \(Cx=0\) iff \(ABx=0\); equal kernels on the \(m\)-dimensional
domain give equal ranks, proving (a). The rank is zero exactly when
\(AB=0\), giving (b).
If \(A\succ0\), multiplication by \(A^{1/2}\) preserves rank, so
the Gram matrix has rank \(\operatorname{rank}B\).
It is PD exactly when \(B\) has rank \(m\), and zero exactly when \(B=0\).
For (d), take \(A=\operatorname{diag}(1,0)\) and \(B=(1,0)^\top\).
Then \(B^\top AB=[1]\succ0\), but \(A\) is singular.
\end{solution}

\begin{problem}[Congruence preserves the sign of a form]\label{ex:09:congruence}
If \(A\) is real symmetric and \(B\) is invertible of the same order,
prove
\[
 B^\top AB\succ0\iff A\succ0,\qquad
 B^\top AB\succeq0\iff A\succeq0.
\]
\end{problem}
\begin{solution}
For every \(x\),
\(x^\top B^\top ABx=(Bx)^\top A(Bx)\).
An invertible \(B\) maps all vectors bijectively onto all vectors and
nonzero vectors onto nonzero vectors. The required strict or weak
inequality therefore holds on one side exactly when it holds on the
other. The two matrices need not have the same eigenvalues; this is
a statement about signs of quadratic values.
\end{solution}

\begin{problem}[Eigenvalues between zero and one]\label{ex:09:shrinkage}
Let \(A\succeq0\), \(B\succ0\). Prove every eigenvalue of
\((A+B)^{-1}A\) lies in \([0,1)\).
\end{problem}
\begin{solution}
Set \(C=A+B\succ0\).
The given matrix is similar, via \(C^{1/2}\), to
\(H=C^{-1/2}AC^{-1/2}\), which is symmetric PSD.
Moreover
\(I-H=C^{-1/2}BC^{-1/2}\succ0\).
For a unit eigenvector of \(H\), these two sign statements give
\(\lambda\geqslant0\) and \(1-\lambda>0\).
Similarity preserves the eigenvalues, proving the result even though
\((A+B)^{-1}A\) need not itself be symmetric.
\end{solution}

\begin{problem}[Principal submatrices]\label{ex:09:principal}
Prove every principal submatrix of a PSD matrix is PSD, and every
principal submatrix of a PD matrix is PD.
\end{problem}
\begin{solution}
For selected coordinates indexed by \(S\), let \(A_S\) be the principal
submatrix. Extend a vector \(y\) on those coordinates to \(x\) by
putting zeros in all other coordinates. Then \(y^\top A_Sy=x^\top Ax\).
Nonnegativity follows for every \(y\), and \(y\ne0\) gives \(x\ne0\),
so strict positivity is also inherited.
\end{solution}

\begin{problem}[*Positive leading principal minors]\label{ex:09:leading-minors}
Prove a real symmetric \(A\) is PD if and only if all leading principal
minors are positive.
\end{problem}
\begin{solution}
If \(A\) is PD, every leading principal submatrix is PD and has positive
determinant, proving necessity.
For sufficiency induct on \(n\). The case \(n=1\) is immediate.
Write \(A=\begin{pmatrix}a&b^\top\\b&D\end{pmatrix}\), with \(a>0\).
For each leading principal submatrix \(S_k\) of
\(S=D-bb^\top/a\), the Schur determinant formula gives
\(\det S_k=\det A_{k+1}/a>0\).
By induction \(S\succ0\).
For \(x=(t,y^\top)^\top\),
\[
 x^\top Ax=a\left(t+\frac{b^\top y}{a}\right)^2+y^\top Sy.
\]
This is strictly positive if \(y\ne0\), or if \(y=0,t\ne0\).
Hence \(A\succ0\).
\end{solution}

\begin{problem}[*All principal minors for semidefiniteness]\label{ex:09:all-minors}
Prove a real symmetric \(A\) is PSD if and only if all its principal
minors, not merely the leading ones, are nonnegative.
\end{problem}
\begin{solution}
Necessity follows because principal submatrices of a PSD matrix are PSD
and their eigenvalue products are nonnegative.
Conversely expand \(\det(tI+A)\) by multilinearity in the columns.
Selecting columns from \(A\) at positions in \(S\), and \(te_j\)
elsewhere, contributes \(t^{n-|S|}\det A_S\).
The cofactor sign for selecting the same omitted rows and columns is
positive. Therefore
\[
 \det(tI+A)=\sum_{S\subseteq\{1,\ldots,n\}}t^{n-|S|}\det A_S,
 \qquad \det A_\varnothing=1.
\]
If all principal minors are nonnegative, this expression is strictly
positive for \(t>0\), because it includes \(t^n\).
A negative eigenvalue \(\lambda\) of \(A\) would make it zero at
\(t=-\lambda>0\), a contradiction.
All eigenvalues are nonnegative, so \(A\succeq0\).
\end{solution}

\begin{problem}[Small minors and one zero determinant]\label{ex:09:small-minors}
\begin{enumerate}
\item For real \(a,b,p\), classify
\[
 A=\begin{pmatrix}0&0\\0&p\end{pmatrix},\qquad
 B=\begin{pmatrix}1&a&b\\a&a^2&ab\\b&ab&p\end{pmatrix}
\]
as PD or PSD, and list their principal minors. Explain why their
leading principal minors are nonnegative for every \(p\).
\item If a symmetric \(n\times n\) matrix has positive leading principal
minors of orders \(1,\ldots,n-1\), but determinant zero, prove it is PSD.
\end{enumerate}
\end{problem}
\begin{solution}
Neither displayed matrix is PD: \(Ae_1=0\), and
\(B(a,-1,0)^\top=0\).
For \(A\), the principal minors are \(0,p\) and \(0\), so it is PSD
exactly when \(p\geqslant0\).
For \(B\), the order-one minors are \(1,a^2,p\); the order-two minors,
on indices \(12,13,23\), are
\(0,p-b^2,a^2(p-b^2)\), and its determinant is zero.
Its quadratic form is
\((x+ay+bz)^2+(p-b^2)z^2\), hence it is PSD exactly when \(p\geqslant b^2\).
The leading minors are \(0,0\) for \(A\), and \(1,0,0\) for \(B\),
regardless of \(p\).

For (b), the case \(n=1\) is the zero scalar matrix.
For \(n>1\), partition as
\(M=\begin{pmatrix}C&b\\b^\top&d\end{pmatrix}\), where \(C\succ0\)
by the leading-minor criterion. The determinant formula gives
\(0=\det C(d-b^\top C^{-1}b)\), so \(d=b^\top C^{-1}b\).
Then
\[
 (x^\top,t)M(x^\top,t)^\top
 =(x+tC^{-1}b)^\top C(x+tC^{-1}b)\geqslant0.
\]
\end{solution}

\begin{problem}[Adding a PSD matrix increases determinant]\label{ex:09:det-monotone}
If \(A\succ0\), \(B\succeq0\), prove
\(\det(A+B)\geqslant\det A\), with equality exactly when \(B=0\).
\end{problem}
\begin{solution}
Set \(C=A^{-1/2}BA^{-1/2}\succeq0\).
Factoring gives
\(\det(A+B)=\det A\,\det(I+C)
=\det A\prod_i(1+\mu_i)\), where \(\mu_i\geqslant0\).
Each factor is at least one. Equality holds exactly when all \(\mu_i=0\),
which by spectral decomposition means \(C=0\), equivalently \(B=0\).
\end{solution}

\begin{problem}[Real eigenvalues of a nonsymmetric product]\label{ex:09:real-product}
If \(A,B\) are real symmetric and at least one is PSD, prove that
\(p_{AB}\) splits into real linear factors.
\end{problem}
\begin{solution}
Suppose \(A\succeq0\).
Set \(X=A^{1/2}\) and \(Y=A^{1/2}B\).
Then \(XY=AB\), while \(YX=A^{1/2}BA^{1/2}\) is symmetric.
The equal-size product identity gives identical characteristic
polynomials for \(XY,YX\), including their zero roots.
The real spectral theorem factors the characteristic polynomial of
\(YX\) into real linear factors. The same is therefore true for
\(AB\), even when \(A\) is singular.
If \(B\succeq0\), apply the argument to \(BA\), which has the same
characteristic polynomial as \(AB\).
\end{solution}

\begin{problem}[Simultaneous congruence with one positive form]\label{ex:09:simultaneous-congruence}
Let \(A\succ0\) and \(B=B^\top\) have the same order.
Prove there are invertible \(T\) and real diagonal \(\Lambda\) such that
\[
 A=TT^\top,\qquad B=T\Lambda T^\top,
\]
where \(\Lambda\) contains the eigenvalues of \(A^{-1}B\).
\end{problem}
\begin{solution}
The matrix \(H=A^{-1/2}BA^{-1/2}\) is real symmetric.
Choose orthogonal \(Q\) with \(H=Q\Lambda Q^\top\), and put
\(T=A^{1/2}Q\).
Then \(TT^\top=A\) and \(T\Lambda T^\top=B\).
Also
\(A^{1/2}(A^{-1}B)A^{-1/2}=H\), so \(A^{-1}B\) is similar to \(H\)
and has the eigenvalues on the diagonal of \(\Lambda\).
\end{solution}

\begin{problem}[Subtract a rank-one matrix]\label{ex:09:downdate}
For \(A\succ0\) and a real column \(a\), prove
\(A-aa^\top\succ0\) if and only if \(a^\top A^{-1}a<1\).
\end{problem}
\begin{solution}
Congruence by \(A^{-1/2}\) turns the matrix into \(I-uu^\top\),
where \(u=A^{-1/2}a\).
If \(u=0\), this is \(I\) and the condition is \(0<1\).
Otherwise it acts as \(1-\norm u^2\) on \(\operatorname{span}u\)
and as \(1\) on \(u^\perp\).
Thus it is PD exactly when \(\norm u^2=a^\top A^{-1}a<1\).
\end{solution}

\begin{problem}[Diagonal blocks and their order]\label{ex:09:block-sign}
\begin{enumerate}
\item Prove \(\operatorname{diag}(A,D)\) is PD (respectively PSD)
exactly when both \(A,D\) are PD (respectively PSD).
\item Prove
\[
 \begin{pmatrix}A&B\\B^\top&D\end{pmatrix}\succ0
 \quad\Longleftrightarrow\quad
 \begin{pmatrix}D&B^\top\\B&A\end{pmatrix}\succ0.
\]
\end{enumerate}
\end{problem}
\begin{solution}
For (a), the quadratic form is \(x^\top Ax+y^\top Dy\).
Both positive terms give the stated sign; necessity follows by setting
one of \(x,y\) to zero.
For (b), interchange the coordinate groups \(x,y\).
This is an invertible permutation of all nonzero vectors, and the
quadratic value is unchanged. Equivalently the matrices are congruent
by an orthogonal block permutation.
\end{solution}

\begin{problem}[Fischer's determinant inequality]\label{ex:09:fischer}
If \(Z=\begin{pmatrix}A&B\\B^\top&D\end{pmatrix}\succeq0\), prove
\(\det Z\leqslant\det A\det D\).
\end{problem}
\begin{solution}
First suppose \(Z\succ0\). Its principal block \(A\) is PD, and testing
vectors \((-A^{-1}By,y^\top)^\top\) shows
\(S=D-B^\top A^{-1}B\succ0\).
The determinant formula gives \(\det Z=\det A\det S\).
Since \(D=S+B^\top A^{-1}B\) and the second term is PSD,
the determinant monotonicity result gives \(\det D\geqslant\det S\).
Multiplying by \(\det A>0\) proves the inequality.

For general PSD \(Z\), apply this to \(Z+\epsilon I\), \(\epsilon>0\).
Its diagonal blocks are \(A+\epsilon I,D+\epsilon I\), giving
\[
 \det(Z+\epsilon I)
 \leqslant\det(A+\epsilon I)\det(D+\epsilon I).
\]
Each determinant is a polynomial in \(\epsilon\) and hence continuous.
Let \(\epsilon\) decrease to zero to obtain the stated inequality,
including singular diagonal blocks.
\end{solution}

\begin{problem}[Positive Schur complements]\label{ex:09:positive-schur}
For symmetric diagonal blocks in
\(Z=\begin{pmatrix}A&B\\B^\top&D\end{pmatrix}\), prove that \(Z\succ0\)
exactly when \(A\succ0\) and \(D-B^\top A^{-1}B\succ0\).
Prove the equivalent criterion using \(D\) and \(A-BD^{-1}B^\top\).
\end{problem}
\begin{solution}
When \(A\succ0\), complete the square:
\[
 (x^\top,y^\top)Z(x^\top,y^\top)^\top
 =(x+A^{-1}By)^\top A(x+A^{-1}By)
   +y^\top(D-B^\top A^{-1}B)y.
\]
If both named blocks are PD, the right side is positive for every
nonzero pair \((x,y)\): either \(y\ne0\), or \(y=0,x\ne0\).
Conversely \(Z\succ0\) implies \(A\succ0\) by principal submatrices,
and choosing \(x=-A^{-1}By\), \(y\ne0\), makes the second form strictly
positive. Interchanging the two coordinate groups gives the criterion
using \(D\).
\end{solution}

\begin{problem}[*Positive block matrices and contractions]\label{ex:09:contractions}
For real symmetric \(A\) of order \(m\), \(D\) of order \(n\), and
\(B\in\R^{m\times n}\), prove that
\[
 \begin{pmatrix}A&B\\B^\top&D\end{pmatrix}\succeq0
\]
if and only if \(A,D\succeq0\) and there is \(C\in\R^{m\times n}\) with
\(B=A^{1/2}CD^{1/2}\) and \(C^\top C\preceq I_n\).
\end{problem}
\begin{solution}
For sufficiency, put \(u=A^{1/2}x\), \(v=D^{1/2}y\).
The block quadratic form is
\[
 \norm u^2+2u^\top Cv+\norm v^2
 =\norm{u+Cv}^2+v^\top(I-C^\top C)v\geqslant0.
\]

For necessity, the principal blocks are PSD.
If \(x\in\ker A\), testing \((tx,y)\) for all real \(t\) gives
\(2t\,x^\top By+y^\top Dy\geqslant0\), hence \(B^\top x=0\).
Likewise \(y\in\ker D\) implies \(By=0\).
Take orthonormal matrices \(U,V\) spanning the positive eigenspaces of
\(A,D\), and write
\(A=U\Delta U^\top\), \(D=V\Gamma V^\top\), where the diagonal
entries of \(\Delta,\Gamma\) are positive.
The kernel relations give
\(B=UU^\top B VV^\top\).
If either positive eigenspace is zero, then \(B=0\), and \(C=0\) works.

Otherwise define
\(C_0=\Delta^{-1/2}U^\top BV\Gamma^{-1/2}\).
Restrict the original quadratic form to these eigenspaces and scale
coordinates by the inverse square roots. This gives
\(\begin{pmatrix}I&C_0\\C_0^\top&I\end{pmatrix}\succeq0\).
Testing \((-C_0z,z)\) yields
\(z^\top(I-C_0^\top C_0)z\geqslant0\).
Set \(C=UC_0V^\top\).
Then \(A^{1/2}CD^{1/2}=B\), and
\[
 \norm{Cy}^2=\norm{C_0V^\top y}^2
 \leqslant\norm{V^\top y}^2\leqslant\norm y^2.
\]
Thus \(C^\top C\preceq I_n\). No inverse of a singular diagonal block
has been used.
\end{solution}

\begin{problem}[*A bordered Gram matrix]\label{ex:09:bordered-kernel}
Let \(A\succeq0\) have order \(m\), \(B\in\R^{m\times n}\), and
\[
 Z=\begin{pmatrix}A&B\\B^\top&0\end{pmatrix},\qquad G=A+BB^\top.
\]
Give an example where \(A\) is singular and \(Z\) is invertible, and
state whether \(G\) is singular in that example.
Prove that \(Z\) is invertible exactly when \(B\) has full column rank
and \(G\succ0\).
\end{problem}
\begin{solution}
Take scalar blocks \(A=0,B=1\). Then \(Z\) exchanges the two coordinates
and is invertible, while \(G=1\) is also invertible.

In general \(G\succeq0\), and
\(u^\top Gu=u^\top Au+\norm{B^\top u}^2\).
For a PSD matrix, \(u^\top Au=0\) iff \(Au=0\), by its spectral
representation. Therefore
\(\ker G=\ker A\cap\ker B^\top\).
The equations \(Z(u,v)^\top=0\) are
\(Au+Bv=0\), \(B^\top u=0\).
Multiplying the first by \(u^\top\) gives \(u^\top Au=0\), hence
\(Au=0\) and \(Bv=0\).
Thus
\[
 \ker Z=(\ker G)\times(\ker B).
\]
It is zero precisely when \(G\) is invertible, equivalently PD,
and \(B\) has full column rank.
The same kernel decomposition gives the stronger rank formula
\(\operatorname{rank}Z=\operatorname{rank}G+\operatorname{rank}B\).
\end{solution}

\begin{problem}[Inverse and determinant of the bordered matrix]\label{ex:09:bordered-formulas}
Use \(A,B,Z,G\) as in the preceding problem.
\begin{enumerate}
\item If \(Z\) is invertible, put \(H=B^\top G^{-1}B\).
Prove \(H\succ0\) and
\[
 Z^{-1}=
 \begin{pmatrix}
 G^{-1}-G^{-1}BH^{-1}B^\top G^{-1}&G^{-1}BH^{-1}\\
 H^{-1}B^\top G^{-1}&I-H^{-1}
 \end{pmatrix}.
\]
\item Assuming only \(G\) is invertible, prove
\(\det Z=(-1)^n\det G\det(B^\top G^{-1}B)\).
\item If \(A\) is invertible, prove
\(\det Z=(-1)^n\det A\det(B^\top A^{-1}B)\).
Deduce, for PD \(A\), the equality of the two products without
the factor \((-1)^n\).
\end{enumerate}
\end{problem}
\begin{solution}
For (a), invertibility of \(Z\) gives \(G\succ0\) and full column rank
of \(B\), so \(H\succ0\) by the matrix quadratic-form result.
Write the proposed blocks as \(M,N,N^\top,K\), in order.
Then \(B^\top M=0\), \(B^\top N=I\).
Using \(A=G-BB^\top\),
\[
 AM=GM=I-BN^\top,\qquad
 AN+BK=B H^{-1}-B+B-BH^{-1}=0.
\]
Thus multiplication by \(Z\) on the left gives the identity.
Both matrices are symmetric, so transposing gives the other product
as well. This proves the inverse formula.

For (b), add \(B\) times the second block row of \(Z\) to its first.
This leaves the determinant unchanged and gives
\(\begin{pmatrix}G&B\\B^\top&0\end{pmatrix}\).
Its Schur complement is \(-B^\top G^{-1}B\), proving (b).
No invertibility of that Schur complement is required.
If \(A\) is invertible, apply the same Schur formula directly to \(Z\),
giving (c). For \(A\succ0\), \(G=A+BB^\top\succ0\), so both formulas
apply to the same determinant. Cancelling \((-1)^n\) gives
\[
 \det G\det(B^\top G^{-1}B)=\det A\det(B^\top A^{-1}B),
\]
even if \(B\) is not full column rank and both sides vanish.
\end{solution}

\subsection*{Idempotents and projections}
In this section \(P^2=P\) means that \(P\) is idempotent. Matrices are
real; symmetry is imposed only when stated.

\begin{problem}[Basic idempotent operations]\label{ex:09:idempotent-basics}
\begin{enumerate}
\item Determine the diagonal entries of a diagonal idempotent.
\item If \(P\) is idempotent, show \(P^\top,P^k\) for \(k\geqslant1\),
and \(I-P\) are idempotent. When is \(-P\) idempotent?
\item Show \(\operatorname{diag}(P,Q)\) is idempotent exactly when both
diagonal blocks are idempotent.
\end{enumerate}
\end{problem}
\begin{solution}
A diagonal entry satisfies \(d^2=d\), so it is \(0\) or \(1\).
Transposing \(P^2=P\) gives \((P^\top)^2=P^\top\).
Induction gives \(P^k=P\), while
\((I-P)^2=I-2P+P^2=I-P\).
For \(-P\), the required identity is \(P=-P\), equivalent over the
reals to \(P=0\).
Finally block multiplication gives
\(\operatorname{diag}(P,Q)^2=\operatorname{diag}(P^2,Q^2)\),
proving the last equivalence.
\end{solution}

\begin{problem}[Nonsymmetric examples and a false converse]\label{ex:09:idempotent-examples}
\leavevmode\par
\begin{enumerate}
\item Construct a nonsymmetric idempotent of rank greater than one.
\item Prove an idempotent's eigenvalues belong to \(\{0,1\}\).
Does an eigenvalue list contained in \(\{0,1\}\) imply idempotence?
\end{enumerate}
\end{problem}
\begin{solution}
Take \(P=\begin{pmatrix}1&0&1\\0&1&0\\0&0&0\end{pmatrix}\).
Its square is itself, its first two columns are independent and its
image is their span, so its rank is two; it is not symmetric.
If \(Pv=\lambda v\), then \(Pv=P^2v=\lambda^2v\), giving
\(\lambda^2=\lambda\).
Conversely \(\begin{pmatrix}1&1\\0&1\end{pmatrix}\) has only
eigenvalue \(1\), but its square is
\(\begin{pmatrix}1&2\\0&1\end{pmatrix}\), so it is not idempotent.
The missing property is diagonalizability, which idempotence itself
does imply.
\end{solution}

\begin{problem}[Ordering orthogonal projectors]\label{ex:09:projector-order}
If \(P,Q\) are symmetric idempotents of the same order, prove
\(P\succeq Q\) if and only if \(PQ=Q\).
\end{problem}
\begin{solution}
Suppose \(P-Q\succeq0\). If \(u\in\operatorname{im}Q\), then \(Qu=u\),
so
\(0\leqslant u^\top(P-Q)u=\norm{Pu}^2-\norm u^2\leqslant0\).
The last inequality is the Pythagorean decomposition
\(u=Pu+(I-P)u\) for an orthogonal projector.
Equality forces \((I-P)u=0\); hence \(Pu=u\) on \(\operatorname{im}Q\),
giving \(PQ=Q\).
Conversely, \(PQ=Q\) and transposition give \(QP=Q\). Thus
\((P-Q)^2=P-Q\), and \(P-Q\) is symmetric idempotent.
Its eigenvalues are \(0,1\), so it is PSD.
\end{solution}

\begin{problem}[Diagonal form and trace of an idempotent]\label{ex:09:idempotent-diagonal}
Let \(P\) be an idempotent of order \(n\).
\begin{enumerate}
\item Prove \(\R^n=\operatorname{im}P\oplus\ker P\), and construct a
similarity \(T^{-1}PT=\operatorname{diag}(I_r,0)\).
\item Construct an orthogonal, matrix \(S\) with
\(S^\top PS=\begin{pmatrix}I_r&C\\0&0\end{pmatrix}\).
\item Prove \(\operatorname{rank}P=\tr P=r\), where \(r\) is the
algebraic multiplicity of eigenvalue \(1\).
\item Prove rank \(n\) is equivalent to \(P=I\), and all eigenvalues
zero is equivalent to \(P=0\).
\end{enumerate}
\end{problem}
\begin{solution}
Every vector decomposes as
\(x=Px+(x-Px)\), with first term in the image and second in the kernel.
If \(y=Pz\) and \(Py=0\), then \(y=P^2z=Py=0\), proving directness.
Take \(T\)'s first columns as a basis of the image, on which \(P\)
acts as the identity, and the remaining columns as a basis of the kernel,
on which it acts as zero. This proves (a), with \(r=\operatorname{rank}P\).

For (b), choose instead an orthonormal basis of the image and extend it
to an orthonormal basis of the whole space. The first columns are fixed
by \(P\), and every remaining column is mapped into the image.
Therefore its matrix has exactly the displayed form.
The diagonalization in (a) gives characteristic polynomial
\((t-1)^r t^{n-r}\) and trace \(r\), since trace is similarity invariant.
This proves (c).
If rank is \(n\), the image is the whole space and \(P\) acts as the
identity. If all eigenvalues are zero, \(r=0\), so the rank is zero
and \(P=0\). Both converses are immediate.
\end{solution}

\begin{problem}[A rank characterization of idempotence]\label{ex:09:idempotent-ranks}
For any real square \(A\) of order \(n\), prove
\[
 A^2=A\quad\Longleftrightarrow\quad
 \operatorname{rank}A+\operatorname{rank}(I-A)=n.
\]
In particular this holds for symmetric \(A\).
\end{problem}
\begin{solution}
If \(A^2=A\), then
\(\operatorname{im}(I-A)=\ker A\), by the decomposition in the
preceding problem. Rank--nullity gives the rank equality.
Conversely \(x=Ax+(I-A)x\) shows
\(\operatorname{im}A+\operatorname{im}(I-A)=\R^n\).
If the dimensions sum to \(n\), the intersection is zero.
But for every \(x\),
\(A(I-A)x=(I-A)Ax\) belongs to both images, so it is zero.
Hence \(A-A^2=0\).
\end{solution}

\begin{problem}[When a sum is idempotent]\label{ex:09:idempotent-sum}
Let \(P,Q\) be idempotents of the same order.
\begin{enumerate}
\item Prove \(P+Q\) is idempotent exactly when \(PQ=QP=0\).
\item In that case, prove
\(\operatorname{rank}(P+Q)=\operatorname{rank}P+\operatorname{rank}Q\).
\end{enumerate}
\end{problem}
\begin{solution}
Expanding gives \((P+Q)^2=P+Q+PQ+QP\).
If the sum is idempotent, \(PQ+QP=0\).
Multiplying this equation by \(P\) on the left and on the right gives
\(PQ+PQP=0\) and \(PQP+QP=0\), respectively.
Subtracting yields \(PQ=QP\); combined with their zero sum, this gives
\(PQ=QP=0\). The converse follows from the expansion.
All three matrices are now idempotent, so rank equals trace for each:
\[
 \operatorname{rank}(P+Q)=\tr(P+Q)=\tr P+\tr Q
 =\operatorname{rank}P+\operatorname{rank}Q.
\]
No symmetry was used.
\end{solution}

\begin{problem}[Two absorption identities]\label{ex:09:absorption}
If \(AB=A\) and \(BA=B\), prove that both \(A,B\) are idempotent.
\end{problem}
\begin{solution}
Associativity gives
\(A^2=(AB)A=A(BA)=AB=A\).
Likewise \(B^2=(BA)B=B(AB)=BA=B\).
\end{solution}

\begin{problem}[Two complementary column spaces]\label{ex:09:complementary}
Let \(A=(A_1\ A_2)\) be square of order \(n\), with
\(\operatorname{rank}A_1+\operatorname{rank}A_2=n\) and
\(A_1^\top A_2=0\). Prove
\[
 A_1(A_1^\top A_1)^{-1}A_1^\top
 +A_2(A_2^\top A_2)^{-1}A_2^\top=I_n.
\]
\end{problem}
\begin{solution}
If the respective column counts are \(k,n-k\), their ranks are at most
these counts; the rank sum forces equality in each. Thus both Gram
matrices are invertible.
The column spaces are orthogonal by \(A_1^\top A_2=0\), and their
dimensions sum to \(n\); hence their orthogonal direct sum is \(\R^n\).
Each displayed term is the orthogonal projector onto its corresponding
column space. Their sum therefore fixes every vector and equals \(I_n\).
The projector formula follows directly by solving the normal equations
for projection, or by an orthonormal basis of each column space.
Empty column groups may be omitted from the formula.
\end{solution}

\subsection*{Extrema and complete numerical examples}

\begin{problem}[Convex and concave extremal eigenvalues]\label{ex:09:eigen-convexity}
For symmetric \(A,B\) of order \(n\) and \(0\leqslant\alpha\leqslant1\),
prove
\[
 \lambda_{\max}(\alpha A+(1-\alpha)B)
 \leqslant\alpha\lambda_{\max}(A)+(1-\alpha)\lambda_{\max}(B),
\]
\[
 \lambda_{\min}(\alpha A+(1-\alpha)B)
 \geqslant\alpha\lambda_{\min}(A)+(1-\alpha)\lambda_{\min}(B).
\]
\end{problem}
\begin{solution}
For every unit \(x\), the Rayleigh bounds give
\[
 \alpha\lambda_{\min}(A)+(1-\alpha)\lambda_{\min}(B)
 \leqslant x^\top\{\alpha A+(1-\alpha)B\}x
 \leqslant\alpha\lambda_{\max}(A)+(1-\alpha)\lambda_{\max}(B).
\]
Taking the minimum in the first inequality and maximum in the second
proves the two claims. The nonnegative weights are essential.
Thus the largest eigenvalue is a convex function of the symmetric
matrix, and the smallest is concave.
\end{solution}

\begin{problem}[Extrema on a circle]\label{ex:09:circle}
Find the eigenvalues and maximum of the quadratic form on the unit
circle for
\[
 A=\begin{pmatrix}2&-1\\-1&2\end{pmatrix},\qquad
 B=\begin{pmatrix}1&1\\1&0\end{pmatrix}.
\]
\end{problem}
\begin{solution}
\(p_A(t)=(t-2)^2-1=(t-1)(t-3)\).
Unit eigenvectors are \((1,1)^\top/\sqrt2\) for \(1\) and
\((1,-1)^\top/\sqrt2\) for \(3\).
The maximum is \(3\), attained at both signs of the latter vector.

For \(B\), \(p_B(t)=t(t-1)-1=t^2-t-1\), giving
\(\lambda_\pm=(1\pm\sqrt5)/2\).
A vector for \(\lambda\) is \((\lambda,1)^\top\), because
\((\lambda+1,\lambda)^\top=(\lambda^2,\lambda)^\top\).
The maximum is \(\lambda_+\), attained at
\(\pm(\lambda_+,1)^\top/\sqrt{\lambda_+^2+1}\).
\end{solution}

\begin{problem}[Extrema on a sphere]\label{ex:09:sphere-example}
Find the eigenvalues and maximum of the associated quadratic form on
the unit sphere for
\(T=\begin{pmatrix}2&-1&0\\-1&2&-1\\0&-1&2\end{pmatrix}\).
\end{problem}
\begin{solution}
Expansion gives
\[
 p_T(t)=(t-2)\{(t-2)^2-1\}-(t-2)
       =(t-2)\{(t-2)^2-2\}.
\]
The eigenvalues are \(2-\sqrt2,2,2+\sqrt2\), with unit eigenvectors
\[
 \tfrac12(1,\sqrt2,1)^\top,\qquad
 \tfrac1{\sqrt2}(1,0,-1)^\top,\qquad
 \tfrac12(1,-\sqrt2,1)^\top,
\]
respectively. Direct multiplication verifies the equations; their
pairwise dot products are zero.
The maximum is \(2+\sqrt2\), attained at both signs of the last vector.
\end{solution}

\begin{problem}[A form with positive and negative extrema]\label{ex:09:mixed-form}
Find the maximum and minimum of \(3x^2+5xy-4y^2\) on \(x^2+y^2=1\).
\end{problem}
\begin{solution}
The symmetric matrix is
\(A=\begin{pmatrix}3&5/2\\5/2&-4\end{pmatrix}\), with
\(p_A(t)=t^2+t-73/4\).
Thus the extrema are the eigenvalues
\[
 \lambda_{\max}=\frac{-1+\sqrt{74}}2,\qquad
 \lambda_{\min}=\frac{-1-\sqrt{74}}2 .
\]
Eigenvectors for these values are
\(v_\pm=(5,-7\pm\sqrt{74})^\top\):
the first equation \(3x+(5/2)y=\lambda x\) holds directly, and the
second follows from the characteristic equation.
Normalize each \(v_\pm\); the maximum and minimum are attained at both
signs of the respective unit vectors. Their dot product is
\(25+49-74=0\), as symmetry predicts.
\end{solution}

\begin{problem}[Symmetric idempotents and their factors]\label{ex:09:projector-factor}
\begin{enumerate}
\item Prove that a real symmetric matrix whose eigenvalues all belong
to \(\{0,1\}\) is idempotent.
\item If \(P\) is a symmetric idempotent of rank \(r\), prove \(P\succeq0\)
and construct an \(n\times r\) matrix \(S\) with
\(S^\top S=I_r\), \(P=SS^\top\).
\end{enumerate}
\end{problem}
\begin{solution}
The spectral theorem gives \(P=Q\Lambda Q^\top\).
For eigenvalues \(0,1\), \(\Lambda^2=\Lambda\), so \(P^2=P\).
Conversely an idempotent's eigenvalues satisfy \(\lambda^2=\lambda\),
so a symmetric idempotent is PSD.
Take \(S\) to consist of the orthonormal eigenvectors for eigenvalue \(1\).
There are \(r\) such vectors, and
\(SS^\top=\sum_{\lambda_j=1}q_jq_j^\top=P\).
For \(r=0\), this is the empty factorization of the zero matrix.
\end{solution}

\subsection*{Functions of matrices: reading practice}
For the following problems, a matrix power series is interpreted by
convergence of every entry. The usual scalar series for the exponential,
logarithm, and binomial function are assumed from calculus wherever they
are invoked.

\begin{problem}[Functions of diagonal matrices]\label{ex:09:diagonal-functions}
\begin{enumerate}
\item If \(D=\operatorname{diag}(d_1,\ldots,d_n)\) and the series converges,
prove
\[
 \sum_{j=0}^\infty c_jD^j
 =\operatorname{diag}\left(\sum_{j=0}^\infty c_jd_1^j,\ldots,
                           \sum_{j=0}^\infty c_jd_n^j\right).
\]
Deduce \(\exp(I_n)=eI_n\).
\item For \(A=\operatorname{diag}(a,b)\), compute \(\exp A\).
If real \(|a|,|b|<1\), compute \((I+A)^\nu\) and \(\log(I+A)\)
from their scalar series.
\item If instead \(|a|,|b|>1\) and \(\nu\) is an integer, compute
\((I+A)^\nu\).
\end{enumerate}
\end{problem}
\begin{solution}
Every finite partial sum is diagonal with the stated scalar partial sums
on its diagonal. Taking their limits proves (a).
Using the exponential series gives \(\exp(I_n)=eI_n\), and
\(\exp A=\operatorname{diag}(e^a,e^b)\).
For \(|a|,|b|<1\), \(1+a,1+b>0\) and the scalar series give
\[
 (I+A)^\nu=\operatorname{diag}((1+a)^\nu,(1+b)^\nu),\qquad
 \log(I+A)=\operatorname{diag}(\log(1+a),\log(1+b)).
\]
In (c), repeated diagonal multiplication gives the same power formula
for nonnegative integers. Since \(|a|,|b|>1\) excludes \(a=b=-1\)
individually, both \(1+a,1+b\) are nonzero; inverse diagonal
multiplication extends the formula to negative integers.
No binomial-series convergence is claimed outside its scalar interval.
\end{solution}

\begin{problem}[Nilpotence terminates a power series]\label{ex:09:nilpotent-functions}
If \(A^k=0\), prove that every formal series
\(\sum_{j=0}^\infty c_jA^j\) reduces to a polynomial of degree at most
\(k-1\).
\end{problem}
\begin{solution}
For \(j\geqslant k\), \(A^j=A^kA^{j-k}=0\).
Thus every partial sum from index \(k-1\) onward equals
\(c_0I+c_1A+\cdots+c_{k-1}A^{k-1}\).
The matrix series therefore requires no convergence assumption on the
coefficients beyond these finitely many terms.
\end{solution}

\begin{problem}[Functions of an idempotent]\label{ex:09:idempotent-functions}
For an idempotent \(P\), prove
\[
 (I+P)^{-1}=I-\tfrac12P,\qquad
 \exp P=I+(e-1)P,\qquad
 \log(I+P)=(\log2)P .
\]
\end{problem}
\begin{solution}
In either multiplication order,
\((I+P)(I-P/2)=I+P/2-P^2/2=I\).
For every \(j\geqslant1\), \(P^j=P\), so the exponential series gives
\[
 \exp P=I+\left(\sum_{j=1}^\infty\frac1{j!}\right)P
       =I+(e-1)P.
\]
The logarithm series has partial sums
\(\sum_{j=1}^N(-1)^{j+1}P^j/j
=\{\sum_{j=1}^N(-1)^{j+1}/j\}P\).
The scalar alternating harmonic series converges to \(\log2\);
hence every matrix entry converges to that of \((\log2)P\).
Symmetry is not needed.
\end{solution}

\begin{problem}[A geometric matrix series]\label{ex:09:geometric-series}
Let \(I-A\) be invertible.
\begin{enumerate}
\item For each integer \(k\geqslant0\), prove
\[
 (I-A)^{-1}=(I-A)^{-1}A^k+\sum_{j=0}^{k-1}A^j,
\]
where the sum is empty for \(k=0\).
\item Assume in addition that \(A\) is symmetric and each of its real
eigenvalues lies strictly between \(-1\) and \(1\).
Prove \((I-A)^{-1}=\sum_{j=0}^\infty A^j\).
\end{enumerate}
\end{problem}
\begin{solution}
The telescoping identity
\((I-A)\sum_{j=0}^{k-1}A^j=I-A^k\)
gives (a) after multiplication by the inverse.
For (b), write \(A=Q\Lambda Q^\top\) using the real spectral theorem.
Then \(A^k=Q\operatorname{diag}(\lambda_1^k,\ldots,\lambda_n^k)Q^\top\).
Every \(\lambda_i^k\) tends to zero, since \(|\lambda_i|<1\).
Each matrix entry is a fixed finite linear combination of these powers,
so \(A^k\) tends entrywise to zero.
Letting \(k\to\infty\) in (a) proves the asserted convergence.
The additional assumptions already imply invertibility of \(I-A\),
since none of its eigenvalues \(1-\lambda_i\) is zero.
\end{solution}

\begin{problem}[Real orthogonal matrices as rotation blocks]\label{ex:09:rotation-blocks}
Let \(R(\theta)=\begin{pmatrix}\cos\theta&-\sin\theta\\
\sin\theta&\cos\theta\end{pmatrix}\).
\begin{enumerate}
\item For orthogonal \(S\), prove
\(S\operatorname{diag}(R(\theta_1),\ldots,R(\theta_k))S^\top\)
is orthogonal with determinant \(1\).
\item Prove the same after appending a final scalar block \(1\).
\item Prove these forms represent every real orthogonal matrix of
determinant \(1\), in even and odd order respectively.
\end{enumerate}
\end{problem}
\begin{solution}
Each rotation block is orthogonal with determinant \(1\).
Block multiplication and orthogonal similarity preserve these facts,
including an appended scalar block \(1\). This proves (a,b).

For (c), let \(A\) be real orthogonal and put \(H=(A+A^\top)/2\).
The matrix \(H\) is real symmetric, so its mutually orthogonal
eigenspaces span \(\R^n\). Since \(AH=HA=(A^2+I)/2\), each eigenspace
\(E_c(H)\) is invariant under \(A\).
For a unit \(u\in E_c(H)\),
\[
 c=u^\top Hu=u^\top Au,\qquad |c|\leqslant\norm u\,\norm{Au}=1.
\]
On \(E_c(H)\), write \(A=cI+K\). The identity \(A+A^\top=2cI\)
shows \(K^\top=-K\) on this space. Orthogonality gives
\[
 I=(cI-K)(cI+K)=c^2I-K^2,\qquad K^2=-(1-c^2)I.
\]
If \(c=\pm1\), then
\(\norm{Ku}^2=-u^\top K^2u=0\), so \(K=0\) and \(A=cI\).
If \(|c|<1\), put \(s=\sqrt{1-c^2}>0\).
For any unit \(u\) in this space, set \(v=Ku/s\). Then
\[
 \norm v=1,\quad u^\top v=0,\quad Ku=sv,\quad Kv=-su.
\]
Thus in the orthonormal basis \((u,v)\) this invariant plane has matrix
\(\begin{pmatrix}c&-s\\s&c\end{pmatrix}=R(\theta)\).
Its orthogonal complement in \(E_c(H)\) is \(K\)-invariant:
for \(z\) perpendicular to the plane and \(w\) in it,
\(w^\top Kz=-(Kw)^\top z=0\).
Repeat on that complement. This constructs all rotation planes using
only real vectors.
The remaining directions carry \(+1\) or \(-1\).
There are an even number of \(-1\) directions because \(\det A=1\);
pair them into \(R(\pi)\). Pair \(+1\) directions into \(R(0)\).
Exactly one \(+1\) direction remains when the total order is odd.
Collecting these orthonormal bases proves (c).
\end{solution}

\begin{problem}[Canonical blocks of a real skew-symmetric matrix]\label{ex:09:skew-blocks}
Let \(A\) be real and \(A^\top=-A\).
\begin{enumerate}
\item Prove that an orthogonal \(S\) exists with
\[
 S^\top AS=\operatorname{diag}
 \left(\begin{pmatrix}0&-\theta_1\\\theta_1&0\end{pmatrix},
 \ldots,\begin{pmatrix}0&-\theta_k\\\theta_k&0\end{pmatrix},0\right),
\]
for real \(\theta_j>0\) and \(k\leqslant n/2\); the zero block may
have size zero.
\item Deduce \(\det A=0\) in odd order and \(\det A\geqslant0\)
in even order.
\end{enumerate}
\end{problem}
\begin{solution}
The matrix \(-A^2=A^\top A\) is real symmetric positive semidefinite.
If it is zero, \(\norm{Ax}^2=x^\top(-A^2)x=0\) for every \(x\),
so \(A=0\).
Otherwise the real spectral theorem supplies a unit \(u\) with
\(-A^2u=\theta^2u\), \(\theta>0\).
Set \(v=Au/\theta\). Then
\[
 \norm v^2=u^\top(-A^2)u/\theta^2=1,\qquad
 u^\top v=u^\top Au/\theta=0,
 \qquad Au=\theta v,\quad Av=-\theta u.
\]
The plane spanned by \(u,v\) is invariant and has the required block.
Its orthogonal complement is invariant too:
\(w^\top Az=-(Aw)^\top z=0\) for \(w\) in the plane and \(z\)
perpendicular to it.
Repeat for the skew-symmetric restriction to that smaller space.
If its squared Gram matrix becomes zero, the restriction itself is zero
by the first argument. Finite induction yields (a).
Every nonzero two-by-two block has determinant \(\theta_j^2>0\).
A positive-size zero block makes the determinant zero.
Odd order necessarily leaves such a block, whereas even order gives
either zero or a product of positive numbers. This proves (b).
\end{solution}

\needspace{22\baselineskip}
\begin{problem}[A power series in a symmetric matrix]\label{ex:09:spectral-functions}
Let \(A=Q\Lambda Q^\top\) be real symmetric with \(Q\) orthogonal.
Suppose the scalar series \(f(t)=\sum_{j=0}^\infty c_jt^j\) converges at
every eigenvalue of \(A\). Prove that the matrix series converges and
\[
 \sum_{j=0}^\infty c_jA^j
 =Q\operatorname{diag}(f(\lambda_1),\ldots,f(\lambda_n))Q^\top.
\]
\end{problem}
\begin{solution}
Cancellation of \(Q^\top Q\) gives \(A^j=Q\Lambda^jQ^\top\) for every
\(j\geqslant0\), including \(j=0\).
Each finite partial sum therefore equals \(Q\) times the diagonal matrix
of the scalar partial sums, times \(Q^\top\).
There are only finitely many diagonal positions, each converging by
assumption. Each entry of the product is a fixed finite linear
combination of them, so it converges to the stated value.
Within a repeated eigenspace the value is the same scalar
\(f(\lambda)\); thus changing its orthonormal eigenbasis does not
change the resulting matrix.
\end{solution}
